\chapter{Summary and Future Plan}
\label{ch:summary}

\section{Summary}

A counterexample-guided repair loop spends a counterexample twice and, in the
usual implementation, wastes it both times: once by running the oracle in an
order nobody chose, and once by letting the proposer redraw from the mode it
has just drawn from. This proposal addresses the first waste and designs for
the second.

The route to it ran through a previous study. \cegmem\ gave the repair loop a
memory keyed by a failure type --- edit location crossed with violated property
--- and that memory worked: oracle calls per episode fell from $11.11$ to
$1.57$, and the repair rate rose from $0.526$ to $0.581$ over $309$ tasks.
Reading its columns carefully is what produced the present question. Oracle
\emph{calls} fell $7.07\times$, but sandbox \emph{runs} fell only $2.35\times$,
because a class-keyed index decides \emph{whether} to pay the oracle and has
nothing to say about \emph{how} the oracle spends once it is paid. A failure
class is a category; a test case is an observation; only the latter can be
ordered, costed, or phrased to a model as a requirement.

\caster\ therefore re-keys the memory by test case and ranks probes by
$K/\hat m^{\,c}$ --- an online Beta-posterior kill rate over an online mean CPU
cost. The preliminary study of Section~\ref{sec:pilot}, run on the frozen
\cegmem\ data because only that data carries the complete patch~$\times$~case
verdict matrix an exact counterfactual replay needs, establishes the following
on $68$ tasks, $2{,}523$ cases and $4{,}670$ patches:

\begin{itemize}[leftmargin=*,itemsep=3pt]
\item \textbf{Against the incumbent}, measured over $6{,}800$ rounds per
policy, \km\ removes $76.2\,\%$ of oracle CPU (95\,\% CI $[-84.9, -61.2]$,
$p = 1.7\times10^{-11}$), and is cheaper on $63$ of $68$ tasks.
\item \textbf{Against the published \frtp\ baseline}, \km\ removes a further
$15.7\,\%$ of CPU (CI $[-28.9, -5.1]$, $p = 0.011$, $42$ of $68$ tasks)
\emph{while paying $23.3\,\%$ more probes}. The mechanism is visible in the
cost per probe: \frtp\ pays 0.375\,s, which is \emph{more} than the
incumbent's 0.336\,s, because ranking on kill count steers it
towards expensive cases; \km\ pays 0.257\,s.
\item \textbf{The unit of account is not neutral.} CPU and executions rank the
two policies oppositely on this corpus, so a prioritisation result reported in
executions --- as the literature reports it, for good reproducibility
reasons --- would have inverted the conclusion.
\item \textbf{The exponent is a parameter, not a constant.} The cost-exponent
sweep has an interior optimum at $c = 0.5$, beating plain \km\ by $9.2\,\%$
(CI $[-15.0,-1.6]$); both $c{=}0$ and $c{=}1.25$ are worse.
\item \textbf{A negative result.} Additional-greedy coverage, the strongest
classical prioritiser, performs barely better than pool order and worse than
random here, because reference coverage is nearly constant across a contest
task's pool (median within-task CV $0.073$). It is dropped from the score.
\item \textbf{Soundness holds.} All three measured policies accepted exactly
the same number of patches, as Definition~\ref{def:soundness} requires.
\end{itemize}

What the pilot does not establish is equally important. It says nothing about
the proposer: RQ3 and RQ4 are untouched, because no replay can manufacture a
counterfactual model call. It is a pilot at $68$ tasks, drawn from a subset
whose selection correlates with cost. And every patch in the matrix came from
one model.

\section{Future Direction}
\label{sec:future}

The remaining work is organised so that the pipeline is frozen before the
evaluation runs, rather than tuned while it runs. \textbf{Every item below is
Phase~2 work and none of it has been executed.}

\begin{itemize}[leftmargin=*,itemsep=4pt]
\item \textbf{Run the Block-A grid on the full corpus.} Thirteen replay arms
over roughly $1{,}244$ gated faults, with the \code{audit} full-scan pass
establishing the kill matrix and the cost table independently of any arm.
This re-establishes RQ1 at scale and runs the ablation ladder of
Section~\ref{sec:ablation} for RQ2.
\item \textbf{Fix $c$ by protocol, not by hindsight.} The pilot's $c = 0.5$
was selected on the pilot arena. Phase~2 declares $c$ before the corpus run and
reports the sweep as sensitivity, so the headline is not a selected maximum.
\item \textbf{Implement and run Block B.} Nine steered arms, each a full
model-call grid. This is the first test of the proposal's second half --- that
a counterexample keyed to a test case and phrased as a positive requirement
moves the proposal distribution where a prohibition and a transcript do not ---
measured against independent repeated sampling, which is the published stronger
strategy~\citep{chen2025conversational}.
\item \textbf{Measure the disclosure dose (RQ4).} The $k$ sweep and the
disclosure sweep are one dose-response curve; \code{km-explore0} must be in
place first or the curve compares $k{=}1$ against itself.
\item \textbf{Report both bills together.} The oracle's CPU and the proposer's
tokens are halves of one cost. \textsc{CigaR} optimised the second and reported
$127$K against $467$K tokens per bug~\citep{hidvegi2024cigar}; this project
optimises the first and should report the pair rather than either alone.
\item \textbf{Beyond this setting (thesis horizon).} Whole-program probes and
single-file faults are what make selection unavailable and prioritisation the
only lever. Carrying the cost-aware ledger into a setting with method-level
suites, where selection \emph{is} available~\citep{lou2024aprregression}, is
the natural generalisation and is out of scope here.
\end{itemize}

\section{Limitations}

\begin{table}[H]
\centering
\caption{Limitations and planned mitigations.}
\label{tab:limitations}
\small
\begin{tabularx}{\textwidth}{@{}X X X@{}}
\toprule
Limitation & Impact & Mitigation\\
\midrule
The \caster\ grid has not been run; evidence is a pilot on the predecessor's data & RQ1 is established at $68$ tasks, not at corpus scale; RQ2 only partially & Phase~2 Block-A grid on $\approx 1{,}244$ faults with the full ablation ladder\\
Arena selection correlates with cost (only tasks whose cost and coverage passes finished) & The $15.7\,\%$ margin over \frtp\ may shrink & The full corpus admits every sound fault; no pre-treatment screen\\
One proposer model (\code{gpt-4o-mini}) produced every patch in the matrix & The kill structure the ledger learns is one model's & A second proposer in Phase~2 sensitivity\\
The replay reuses one cost sample per case & Understates cost variance & The headline rests on the re-executed block, not the replay\\
$c = 0.5$ was selected on the same arena it is reported on & A selected maximum is optimistically biased & Declare $c$ before the corpus run; report the sweep as sensitivity\\
Whole-program, single-file, stdout-compared faults & No claim transfers to project-scale Java repair & Stated as external threat; selection-capable settings are future work\\
CPU is the cost; money and wall clock are not & A differently priced deployment weights the units differently & Both CPU and executions reported for every arm\\
\bottomrule
\end{tabularx}
\end{table}

\section{Timeline}

The project runs across the 2026 calendar year in two phases separated by the
institutional summer break. Phase~1 covers the work reported here: the
\cegmem\ study that generated the question, the literature consolidation, the
\caster\ implementation and its unit tests, and the pilot. Phase~2 covers the
grid, the steered block, and the writing.

\begin{figure}[H]
\centering
\resizebox{\textwidth}{!}{%
\begin{ganttchart}[
  hgrid, vgrid,
  x unit=0.62cm, y unit title=0.58cm, y unit chart=0.52cm,
  title/.style={fill=none,draw=black!55},
  title label font=\scriptsize\bfseries,
  bar/.style={fill=bkblue!50,draw=bkblue,rounded corners=1pt},
  bar height=0.52,
  group/.style={fill=bkblue!85,draw=none},
  group right shift=0, group left shift=0,
  bar label font=\scriptsize,
  group label font=\scriptsize\bfseries,
  milestone label font=\scriptsize\itshape
]{1}{17}
\gantttitle{2026}{12} \gantttitle{2027}{5} \\
\gantttitlelist{"J","F","M","A","M","J","J","A","S","O","N","D","J","F","M","A","M"}{1} \\
\ganttgroup{Phase 1 -- premise and proposal}{1}{11} \\
\ganttbar{\cegmem\ study (typed memory)}{1}{6} \\
\ganttbar{Literature consolidation}{7}{8} \\
\ganttbar{\caster\ design and implementation}{9}{10} \\
\ganttbar{Unit tests and invariants}{9}{10} \\
\ganttbar{Cost / coverage measurement passes}{10}{11} \\
\ganttbar{Pilot: replay + re-execution}{11}{11} \\
\ganttbar{Proposal writing}{11}{11} \\
\ganttmilestone{Proposal defence}{11} \\
\ganttgroup{Phase 2 -- the grid}{12}{17} \\
\ganttbar{Corpus freeze; \code{audit} full scan}{12}{13} \\
\ganttbar{Block-A grid (13 arms)}{13}{14} \\
\ganttbar{Block-B implementation and grid (9 arms)}{14}{15} \\
\ganttbar{Disclosure sweep (RQ4)}{15}{16} \\
\ganttbar{Sensitivity, threats, writing}{16}{17} \\
\ganttmilestone{Final defence}{17}
\end{ganttchart}}
\caption{Project timeline. Everything up to and including the proposal-defence
milestone has been carried out; every bar below it is planned work.
Deliverables per work package are in Table~\ref{tab:workpackages}.}
\label{fig:timeline}
\end{figure}

\begin{table}[H]
\centering
\caption{Work packages by phase.}
\label{tab:workpackages}
\small
\begin{tabularx}{\textwidth}{@{}l l X X@{}}
\toprule
Phase & Period & Main focus & Deliverable\\
\midrule
\multirow{5}{*}{Phase 1} & Jan--Jun & \cegmem: typed memory, the $309$-task grid & The premise, and the residual that generated RQ1\\
 & Jul--Aug & Literature consolidation & Related-work matrix; the gap statement of Sec.~\ref{sec:gap}\\
 & Sep--Oct & \caster\ design and implementation & Ten verbs; the three module boundaries; 263 unit tests\\
 & Oct--Nov & Cost and coverage measurement passes & \code{cost\_all}, \code{cov\_all}, the verdict matrix\\
 & Nov & Pilot and proposal & Sec.~\ref{sec:pilot}; this document; defence\\
\midrule
\multirow{5}{*}{Phase 2} & Dec--Jan & Corpus freeze; \code{audit} full scan & Kill matrix and frozen cost table, arm-independent\\
 & Jan--Feb & Block-A grid, 13 arms & RQ1 at corpus scale; the RQ2 ablation ladder\\
 & Feb--Mar & Block-B implementation and grid, 9 arms & RQ3: duplicate rate against repeated sampling\\
 & Mar--Apr & Disclosure sweep & RQ4 dose-response curve\\
 & Apr--May & Sensitivity, threats, writing & Final manuscript and defence\\
\bottomrule
\end{tabularx}
\end{table}

Two buffers are built into the plan. The evaluation protocol --- the corpus
gate, the pool order, the cost table and the declared value of $c$ --- is
frozen at the end of Phase~1, so unexpected behaviour in Phase~2 can be
analysed rather than hidden by continuous retuning. And because Block~A costs
no model calls, the thirteen replay arms can be re-run freely whenever the
analysis changes, which is what makes the ablation ladder affordable at all.
